\subsection{The five-element middle with four and six adjacent arrows}

\begin{proposition}\label{C:middle-five-four-six}
Let $C$ have objects $0<1<2$, trivial endomorphism monoids at $0$ and
$2$, a five-element middle monoid $N$, and hom-set cardinalities
\[
 |C(0,1)|=4,\qquad |C(1,2)|=6,\qquad |C(0,2)|=1.
\]
Then $BC$ is not homotopy equivalent to $\CP^2$.
The same conclusion holds for the opposite category.
\end{proposition}

\begin{proof}
Write $A=\Q N$, $V=\ker(\Q C(0,1)\to\Q)$ and
$W=\ker(\Q C(1,2)\to\Q)$.  Their dimensions are three and five.
The exact low-degree sequence and high-degree formula of
Proposition~\ref{C:linear-augmentation-Tor}, with its single long arrow,
would give
\[
 \dim\operatorname{Tor}_{0,1,2}^{A}(W,V)=(1,0,1),
 \qquad \operatorname{Tor}_i^A(W,V)=0\quad(i\geq3).
\]
The algebra $\Q N$ is nonlocal.  Indeed, locality would force every
idempotent of the finite monoid $N$ to be its identity, since a distinct
monoid basis element is neither the zero algebra element nor the algebra
identity.  A finite monoid with only that idempotent is a group.
Maschke's theorem then makes $\Q N$ semisimple; a semisimple local algebra
is a division algebra, whose augmentation onto $\Q$ is an isomorphism.
This contradicts $|N|=5$.

Adjoining a new absorbing zero exhibits $A$ as a contracted monoid
algebra, so Proposition~\ref{B:five-nonlocal-exceptions} leaves precisely
the two exceptional algebras.  Use the indecomposable module notation and
Tor tables of Lemma~\ref{B:five-exceptional-targets}.
For the line algebra, the degree-two condition forces exactly one
$S_3$ in $V$ and one $T_1$ in $W$; the degree-one condition excludes
$S_2$ from $V$ and $T_2$ from $W$.  Hence
\[
 V\in\{2S_1\oplus S_3,\ S_3\oplus L_a,\ S_3\oplus L_b\},
 \qquad W=T_1\oplus T_3^s\oplus R_a^p\oplus R_b^q,
 \quad s+2p+2q=4.
\]
In these three cases the tensor dimensions are respectively
$2+s+2p$, $s+p+q$, and $2s+q$.
The first cannot equal one.  In the second, combining a value of one
with the displayed dimension equation would give $s=-2$.
The third gives $s=0$, $p=q=1$.  Thus the only line-algebra pair is
\begin{equation}\label{C:middle-five-four-six-pair}
 V=S_3\oplus L_b,\qquad W=T_1\oplus R_a\oplus R_b.
\end{equation}
For the Nakayama algebra, its table gives exactly two pairs with these
module dimensions and with $(\dim\operatorname{Tor}_1,
\dim\operatorname{Tor}_2)=(0,1)$:
\[
 \begin{array}{c|c|c}
 V&W&\dim(W\otimes_A V)\\\hline
 S_2\oplus Ae_2&R\oplus e_1A&2\\
 S_2\oplus Ae_2&T_2\oplus(e_2A)^2&6.
 \end{array}
\]
Neither has tensor dimension one.  The small exact enumeration of these
displayed tables is retained in
\path{code/b_five_middle_Tor_profiles.py}; it makes no finite-set
realizability assumption.

It remains to examine actual finite-set realizations of
\eqref{C:middle-five-four-six-pair}.  The complete multiplicative-basis
classification in \path{code/outputs/b_five_line_monoid_bases.json}
has 48 oriented contracted line bases.  An ordinary five-element monoid
basis is exactly a retained basis whose multiplication never uses the
contracted zero.  There are eight such records, with indices
$0,1,2,3,30,31,32,33$.  This is complete because adjoining a new zero to
any ordinary five-element middle produces a member of the complete
six-element source census.

For each of these eight literal monoids, enumerate all left actions on
four points and all right actions on six points.  The already audited
pointwise action constraint algorithm enforces every equation
$L_{mn}=L_mL_n$ or $R_{mn}=R_nR_m$, starting with the literal identity
action; it does not assume that abstract modules lift to finite sets.
Let $\chi_1,\chi_2,\chi_3$ be the stored ordered simple characters and
$a,b$ the stored radical-arrow vectors.  The exact selection conditions
on the augmentation modules are
\[
 \begin{array}{c|c|c|c}
 &\operatorname{tr}(m)&\operatorname{rank}(a)&\operatorname{rank}(b)\\\hline
 V&\chi_2(m)+2\chi_3(m)&0&1\\
 W&2\chi_1(m)+2\chi_2(m)+\chi_3(m)&1&1.
 \end{array}
\]
Here the trace is the number of fixed points minus one.  Zero and
rank-one conditions are checked integrally, using all two-by-two minors.
The simple multiplicities and the ranks of the two arrows uniquely
specify the decompositions in \eqref{C:middle-five-four-six-pair}.

The complete action search examines 5,112 labelled four-point actions
and 1,488,464 labelled six-point actions.  State conjugacy gives the
following retained action counts:
\[
 \begin{array}{c|rrrrrrrr}
 \text{basis index}&0&1&2&3&30&31&32&33\\\hline
 \text{left actions}&1&1&1&7&1&7&1&2\\
 \text{right actions}&8&4&28&2&29&2&11&4\\
 \text{pairs}&8&4&28&14&29&14&11&8.
 \end{array}
\]
The two sets may be relabelled independently because their composition
has a one-point target.  Thus these 116 pairs cover every possible
category in the remaining module case.

For all 116 literal categories the normalized bar complex over $\F_2$
has chain dimensions $(15,80,320,1280)$ in degrees one through four,
boundary ranks $(2,13,66,254)$, and therefore
$(b_1,b_2,b_3)=(0,1,0)$.  A nonboundary degree-two cocycle and an explicit
degree-three primitive for its square are retained for every category.
The separate Python verifier rebuilds the normalized complexes and
checks all 70,528 typed associativity equations, 37,120 cocycle equations,
and 148,480 square-primitive equations.  It also checks the full
Cartesian linkage from actual action pairs to the literal tables.
Thus every remaining degree-two generator has zero square over $\F_2$,
contrary to $H^*(\CP^2;\F_2)$.
\end{proof}

\begin{remark}\label{C:middle-five-four-six-certificates}
The source/action wrapper is
\path{code/b_middle5_triangle18_actions.cpp}; the literal categories and
their square primitives are
\path{code/outputs/b_middle5_triangle18_categories.json} and
\path{code/outputs/b_middle5_triangle18_cups.jsonl}.
The independent bar/action-linkage check is
\path{code/b_verify_middle5_triangle18.py}, with result
\path{code/outputs/b_middle5_triangle18_independently_verified.json}.
The latter checks the complete retained output and all exact equations;
it does not claim a second independent numerical action enumeration.
\end{remark}
